% TOP_PROBLEM: {"id": 3159, "title": "Partitioning edge-connectivity", "category_id": 3, "category": {"id": 3, "name": "graph_theory", "display_name": "Graph Theory", "description": "Problems involving graphs, networks, and their properties.", "slug": "graph-theory", "order_index": 3, "created_at": "2026-07-31T15:26:25.671Z"}, "status": "open", "problem_number": "OPG-146", "set_id": 12}
% FIRST_POSTED: 2026-10-01
% ATTEMPT_STATUS: unresolved
% UnsolvedMath ID 3159; source catalogue code OPG-146
% !TeX program = lualatex
% GPT-6 Astra Ultra research attempt; independent partial work, 2026-10-01.
\documentclass[11pt,a4paper]{article}
\usepackage[margin=25mm]{geometry}
\usepackage{fontspec}
\setmainfont{lmroman10-regular.otf}[BoldFont=lmroman10-bold.otf,
 ItalicFont=lmroman10-italic.otf,BoldItalicFont=lmroman10-bolditalic.otf]
\usepackage{amsmath,amssymb,mathtools}
\usepackage{unicode-math}
\setmathfont{latinmodern-math.otf}
\AtBeginDocument{\let\setminus\smallsetminus}
\usepackage{xurl,hyperref}
\hypersetup{colorlinks=true,urlcolor=blue}
\setcounter{secnumdepth}{0}
\setlength{\parindent}{0pt}
\setlength{\parskip}{5pt}
\setlength{\emergencystretch}{3em}
\begin{document}
\section*{Partitioning edge-connectivity}\label{3159}
{\small UnsolvedMath ID 3159 · OPG-146 · Graph Theory · OpenGarden}
\subsection{Definitions and mathematical statement}
Let $a,b\geq1$ be integers and let $G=(V,E)$ be a finite
undirected loopless multigraph on at least two vertices. Parallel
edges are counted separately. For a nonempty proper subset
$X\subset V$, write $\delta_G(X)$ for the set of edges having
exactly one endpoint in $X$. The graph is $k$-edge-connected
when $|\delta_G(X)|\geq k$ for every such $X$; equivalently,
deleting fewer than $k$ edges cannot disconnect it.

The question is whether the hypothesis
\[
 |\delta_G(X)|\geq a+b+2
 \qquad(\varnothing\neq X\subsetneq V)
\]
always permits a partition $E=A\mathbin{\dot\cup}B$ for which
$(V,A)$ is $a$-edge-connected and $(V,B)$ is $b$-edge-connected.
In terms of the same cuts, the desired simultaneous conditions are
\[
 |A\cap\delta_G(X)|\geq a,\qquad
 |B\cap\delta_G(X)|\geq b
 \qquad(\varnothing\neq X\subsetneq V).
\]
The two resulting graphs must be spanning: both retain every
vertex of $G$. Each original edge belongs to exactly one of
them, and the partition has to satisfy every cut constraint
at once. It may depend on $G,a,b$.

This asks for an additive connectivity threshold uniform over
all orders and edge multiplicities. Having enough edges in
each cut separately is only the input condition; the issue is
whether their assignments can be coordinated globally. The
special simple-graph case follows by restricting the allowed
multigraphs.
\subsection{Short English statement}
Can connectivity of at least a+b+2 be split into two spanning edge-disjoint graphs of edge connectivity at least a and b?
\subsection{Research attempt: tree packing and exact rounding on three vertices}
\paragraph{Scope and literature check.}
GPT-6 Astra Ultra, 1 October 2026. Edges are distinct resources even
when parallel, and both output graphs must contain every vertex.
The original problem [S2] asks for an additive threshold. The
Nash-Williams spanning-tree theorem [S3] supplies a stronger general
threshold, while a direct rounding argument below verifies the stated
threshold for graphs with at most three vertices. These are restricted
results, not a claimed new general partition theorem.

\paragraph{1. A certified stronger threshold.}
Put $s=a+b$. Suppose every nontrivial cut has at least $2s$ edges.
For a partition of the vertices into $q\ge2$ nonempty parts, the sum
of the $q$ cut sizes is at least $2sq$. Each edge between different
parts contributes twice, so there are at least $sq\ge s(q-1)$
such edges. For the one-part partition the requirement is zero.
By the spanning-tree packing theorem, $G$ has $s$ edge-disjoint
spanning trees. Assign $a$ of them to $A$ and the other $b$ to $B$.
Put all remaining edges in either class, for example $A$.

Every nontrivial cut meets each spanning tree. Edge disjointness
therefore gives at least $a$ distinct $A$-edges and $b$ distinct
$B$-edges across that cut. This proves the required partition under
$2(a+b)$-edge-connectivity. In particular, when $a=b=1$, this is exactly
the requested threshold four. The theorem is being used as an external
result with its partition hypothesis checked, rather than assumed from
minimum degree alone.

\paragraph{2. Full target threshold for three vertices.}
Let the three edge multiplicities be $m_{12},m_{13},m_{23}$, set
$\lambda=a+b+2$, and suppose each vertex cut has size at least
$\lambda$. Define
\[
 p=\frac{a+1}{\lambda},\qquad
 A_{ij}=\lfloor p m_{ij}\rfloor,\qquad
 B_{ij}=m_{ij}-A_{ij}.
\]
Choose that many actual parallel edges for each class. Any nontrivial
cut in a three-vertex graph consists of exactly two bundles. If its
total multiplicity is $c$, then
\[
 |A\cap\delta(X)|>pc-2\ge a-1,
 \qquad |B\cap\delta(X)|\ge(1-p)c\ge b+1.
\]
The first quantity is an integer, hence is at least $a$; the second
is certainly at least $b$. Complementary vertex cuts are the same
edge sets, so these checks exhaust all cuts. The construction also
allows a bundle of multiplicity zero. On two vertices there is just
one bundle: assign any $a$ edges to $A$ and the remainder to $B$.
Thus all loopless multigraphs of orders two and three satisfy the
conjectured threshold, for every $a,b\ge1$.

\paragraph{3. Why the fractional argument is insufficient in general.}
For arbitrary order, assigning fraction $p$ of each edge to $A$ meets
every fractional cut constraint, with slack: the $A$-mass is at least
$a+1$ and the $B$-mass at least $b+1$. This is not an edge partition.
Rounding each parallel bundle down independently loses less than one
unit per bundle crossing a cut. A three-vertex cut has only two
bundles, which is exactly what made the preceding strict inequality
useful. In a large simple graph, a cut may cross arbitrarily many
one-edge bundles. Indeed independent rounding with $0<p<1$ assigns
zero edges of every such bundle to $A$. Consequently fractional
feasibility plus this rounding rule cannot prove the general claim.

\paragraph{Verification and first gap.}
The companion script [S4] checks the three-bundle formula using exact
integer arithmetic for $1\le a,b\le4$ and all bundle multiplicities
between zero and ten satisfying the input cuts. The proof itself is
uniform and does not rely on that finite range. The unresolved step
is coordinating integral choices across all cuts with only the two
units of slack in the original threshold; the tree-packing argument
uses greater connectivity and does not supply that coordination.

\subsection{Final status}
UNRESOLVED. The conjecture is verified here for order at most three
and for $a=b=1$, and a general stronger threshold is recovered.
Completion estimate: no defensible global percentage.

\subsection{Sources}
\begin{itemize}
\item[S1] ULAM.AI and the UnsolvedMath Contributors, supplied problems.json, record 3159 (OPG-146), OpenGarden collection. Catalogue identity and source transcription; CC BY 4.0.
\url{https://huggingface.co/datasets/ulamai/UnsolvedMath}.
\item[S2] Open Problem Garden, Partitioning edge-connectivity. Original problem and discussion; the mathematical formulation below is expanded from the supplied source record.
\url{http://www.openproblemgarden.org/op/partitioning_edge_connectivity}.
\item[S3] C. St. J. A. Nash-Williams, \emph{Edge-Disjoint Spanning Trees of Finite Graphs}, Journal of the London Mathematical Society 36 (1961), 445--450. Checked 1 October 2026.
\url{https://doi.org/10.1112/jlms/s1-36.1.445}.
\item[S4] GPT-6 Astra Ultra, reproducible exact checks accompanying this batch, 1 October 2026. Repository file \texttt{scripts/check\_attempt\_batch02\_connectivity\_2026\_10\_01.py}; run with Python 3.
\end{itemize}
\end{document}
